\section{Experiments and Validation}

The experiments evaluate the main claim of \SPMPC{}: explicit slosh-state prediction can provide interpretable liquid-response improvement beyond smooth-only control, without unacceptable path-tracking, task-time, or online-computation cost.
The evaluation uses a two-layer evidence chain.
Controlled comparisons are used to separate internal ablation factors and local-planner baselines, while physical trials examine external liquid response and the interpretability of internal model proxies.
Online computation and execution diagnostics are treated separately to verify that the method remains a deployable receding-horizon local planner.

\subsection{Experimental Setting and Metrics}

The experimental scope matches the problem formulation: a standard wheeled mobile base transports an open liquid container along a precomputed geometrically feasible reference path.
Planner comparisons use the same start and goal conditions, reference path, base constraints, task termination condition, and failure criteria.
External baseline planners do not receive liquid-surface state as a control input.
Each trial logs trajectory, velocity commands, control inputs, solve time, internal model proxies, and external liquid observations.
Throughout the evaluation, $H_{\mathrm{model}}^{\mathrm{modal}}$ and $H_{\mathrm{model}}^{\mathrm{full}}$ are interpreted only as internal prediction or diagnostic proxies.
Claims about real liquid motion are based on RGB envelope measurements, liquid-level measurements, or other external observations.

\begin{figure}[t]
  \centering
  \resizebox{\linewidth}{!}{%
  \begin{tikzpicture}[
    font=\footnotesize,
    block/.style={draw=spmpcgray, rounded corners=2pt, line width=0.5pt, fill=spmpclight, minimum width=1.55cm, minimum height=0.58cm, align=center},
    core/.style={draw=spmpcblue, rounded corners=2pt, line width=0.7pt, fill=white, minimum width=1.60cm, minimum height=0.72cm, align=center, font=\footnotesize\bfseries},
    metric/.style={draw=spmpcgray, rounded corners=2pt, line width=0.4pt, fill=white, minimum width=1.15cm, minimum height=0.44cm, align=center, font=\scriptsize},
    arr/.style={-{Latex[length=1.6mm]}, line width=0.5pt, draw=spmpcblue}
  ]
    \node[core] (robot) at (0,0) {Base +\\container};
    \node[block] (planner) at (2.35,0.75) {Planner logs\\traj. / cmd};
    \node[block] (proxy) at (2.35,-0.75) {Model proxy\\slosh};
    \node[block] (camera) at (2.35,1.80) {Surface obs.\\RGB / level};
    \node[core] (eval) at (4.65,0.35) {Metrics};
    \node[metric] (m1) at (6.75,1.45) {liquid};
    \node[metric] (m2) at (6.75,0.65) {proxy};
    \node[metric] (m3) at (6.75,-0.15) {task};
    \node[metric] (m4) at (6.75,-0.95) {runtime};

    \draw[arr] (robot.east) -- (planner.west);
    \draw[arr] (robot.east) -- (proxy.west);
    \draw[arr] (robot.north east) |- (camera.west);
    \draw[arr] (planner.east) -- (eval.west);
    \draw[arr] (proxy.east) -- (eval.west);
    \draw[arr] (camera.east) -| (eval.north);
    \foreach \m in {m1,m2,m3,m4}
      \draw[arr] (eval.east) -- (\m.west);
  \end{tikzpicture}}
  \caption{Experimental measurement pipeline. Trajectories, base commands, solve time, and internal slosh proxies explain planner behavior; real liquid-response claims are supported by external vision or liquid-level observations.}
  \label{fig:experiment_setup}
\end{figure}

\Cref{tab:metrics_plan} summarizes the metrics.
For liquid-response comparisons, high-percentile and RMS metrics are emphasized because single-frame maxima can be affected by isolated visual envelope spikes.
For limited sample sizes, the statistical claim is restricted to representative evidence rather than broad population-level significance.

\begin{table}[t]
  \centering
  \scriptsize
  \setlength{\tabcolsep}{3pt}
  \caption{Evaluation metrics and their roles.}
  \label{tab:metrics_plan}
  \begin{tabular}{p{0.27\linewidth} p{0.35\linewidth} p{0.26\linewidth}}
    \toprule
    Category & Metrics & Purpose \\
    \midrule
    External liquid motion & RGB max-LCR, RMS-LCR, residual oscillation & Test whether external observations support liquid-response improvement \\
    Model diagnostics & Predicted slosh state, $H_{\mathrm{model}}^{\mathrm{modal}}$, $H_{\mathrm{model}}^{\mathrm{full}}$ & Interpret model prediction without replacing external measurements \\
    Model-external consistency & $\gamma_{\mathrm{bias}}$, $\gamma_{\mathrm{abs}}$, RMSE, $\rho$, $\tau^\star$, peak/high-percentile underestimation & Check whether internal proxies explain external observations \\
    Task completion & Success rate, timeout, arrival time & Exclude the explanation that the planner simply stops or slows excessively \\
    Path tracking & Contouring error, path deviation & Verify that the local-planning task remains satisfied \\
    Control smoothness & RMS velocity, acceleration, and angular acceleration & Separate smoothness from slosh-state prediction \\
    Online computation & Solve time, control frequency, solver failures & Verify real-time receding-horizon execution \\
    \bottomrule
  \end{tabular}
\end{table}

To avoid evaluating the method using only internal model quantities, we compare model proxies with an external liquid observation $H_{\mathrm{vis}}(t)$.
Let $H_{\mathrm{model}}^{\mathrm{eval}}(t)$ denote the model proxy used for comparison; it may be either $H_{\mathrm{model}}^{\mathrm{modal}}(t)$ or $H_{\mathrm{model}}^{\mathrm{full}}(t)$, and the selected convention is reported with the results.
Following the practice of comparing low-order model height with video-extracted liquid motion in anti-slosh trajectory planning \cite{C01_Ferrari2026}, we define the signed bias
\begin{equation}
  \gamma_{\mathrm{bias}} =
  100\cdot
  \frac{\int_{t_0}^{t_1}\left(H_{\mathrm{model}}^{\mathrm{eval}}(t)-H_{\mathrm{vis}}(t)\right)\dd t}
  {\int_{t_0}^{t_1}H_{\mathrm{model}}^{\mathrm{eval}}(t)\dd t+\epsilon}.
  \label{eq:gamma_bias}
\end{equation}
The interval $[t_0,t_1]$ spans the valid command window and the post-arrival residual-oscillation window, and $\epsilon$ prevents denominator degeneracy.
A positive $\gamma_{\mathrm{bias}}$ indicates that the model is conservative on average, while a negative value indicates average underestimation of the external metric.
The absolute deviation is
\begin{equation}
  \gamma_{\mathrm{abs}} =
  100\cdot
  \frac{\int_{t_0}^{t_1}\left|H_{\mathrm{model}}^{\mathrm{eval}}(t)-H_{\mathrm{vis}}(t)\right|\dd t}
  {\int_{t_0}^{t_1}H_{\mathrm{vis}}(t)\dd t+\epsilon}.
  \label{eq:gamma_abs}
\end{equation}
Additional diagnostics include RMSE, correlation coefficient $\rho$, phase shift $\tau^\star$, peak error, high-percentile error, and maximum local underestimation.
Model-external consistency is used for interpretation; it is not a substitute for external liquid evidence and does not constitute a formal spill-prevention guarantee.

\subsection{Internal Ablation and Baseline Structure}

The internal ablation tests whether \SPMPC{} is merely a stronger smoothness controller.
All variants share the same path-progress \MPCC{} structure, reference path, and base constraints; they differ only in the slosh model/cost and smoothing weights.
The four variants \Bzero{}, \Bsmooth{}, \Bslosh{}, and \Bours{} are defined in \cref{tab:method_variants}.

\begin{figure}[t]
  \centering
  \resizebox{\linewidth}{!}{%
  \begin{tikzpicture}[
    font=\footnotesize,
    cell/.style={draw=spmpcgray, rounded corners=2pt, line width=0.55pt, fill=white, minimum width=2.15cm, minimum height=0.82cm, align=center},
    ours/.style={draw=spmpcblue, rounded corners=2pt, line width=0.8pt, fill=spmpclight, minimum width=2.15cm, minimum height=0.82cm, align=center, font=\footnotesize\bfseries},
    head/.style={font=\footnotesize\bfseries, text=spmpcblue, align=center},
    row/.style={font=\scriptsize\bfseries, text=spmpcgray, align=right},
    arr/.style={-{Latex[length=1.7mm]}, line width=0.5pt, draw=spmpcblue, shorten >=1pt, shorten <=1pt}
  ]
    \node[head] at (1.45,1.65) {No slosh state};
    \node[head] at (4.35,1.65) {Slosh state};
    \node[row, anchor=east] at (-0.05,0.70) {Weak\\smooth.};
    \node[row, anchor=east] at (-0.05,-0.45) {Strong\\smooth.};

    \node[cell] (b0) at (1.45,0.70) {\Bzero{}\\Baseline};
    \node[cell] (bs) at (4.35,0.70) {\Bslosh{}\\Slosh-aware};
    \node[cell] (bm) at (1.45,-0.45) {\Bsmooth{}\\High smoothing};
    \node[ours] (bo) at (4.35,-0.45) {\Bours{}\\SPMPC};

    \draw[arr] (b0) -- (bs);
    \draw[arr] (b0) -- (bm);
    \draw[arr] (bm) -- (bo);
    \draw[arr] (bs) -- (bo);
  \end{tikzpicture}}
  \caption{Internal ablation structure. The variants form a two-factor decomposition of whether the planner propagates a slosh state and whether stronger control smoothing is applied.}
  \label{fig:ablation_design}
\end{figure}

The comparisons answer distinct questions.
\Bsmooth{} vs.\ \Bzero{} quantifies the effect of stronger smoothness alone.
\Bslosh{} vs.\ \Bzero{} evaluates the independent role of explicit slosh-state prediction.
\Bours{} vs.\ \Bsmooth{} tests whether the state-augmented planner exceeds smooth-only planning.
\Bours{} vs.\ \Bslosh{} characterizes the role of additional smoothing in executable local-command generation.
Only trends supported by external liquid observations are interpreted as real liquid-response improvement; reductions in internal proxies alone are interpreted as model-side reductions in predicted slosh response.

Ordinary local-planner baselines test whether online local planning without liquid-state propagation already covers the task.
Candidate baselines include DWA, TEB, \MpcLocalPlanner{}, and LT-DWA.
They operate at a similar deployment layer and generate executable base commands, but generally do not propagate liquid dynamic memory.
Quantitative comparison requires runnable interfaces, matched base constraints, identical task definitions, and reproducible logs.
Mobile-base anti-slosh reference methods, such as Hamaguchi-style path/speed profiles, input shaping, and offline liquid-constrained trajectory optimization, are included quantitatively only when they can be reproduced under comparable platform and interface conditions; otherwise, they serve as decision-layer references rather than ranked baselines.

The Slosh-risk Reference Governor is evaluated as an incremental module by comparing \Bours{} and \Boursgov{} under the same \SPMPC{} parameters, nominal speed reference, and base constraints.
In addition to liquid, task, and tracking metrics, the governor-specific diagnostics are $\beta_{\mathrm{raw}}$, $\beta_f$, saturated rate, transient rate-limited rate, and arrival-time change.

\subsection{Physical Experiments: External Liquid Response and Model Consistency}

Physical experiments examine whether the trends observed from planning and model proxies are supported by external liquid-surface observations.
The platform records the final base command, external liquid observation, internal model proxy, and execution-chain diagnostics, allowing the analysis to separate planner behavior from base execution, sensing latency, and command limiting.
In the physical analysis, $H_{\mathrm{model}}$ remains an internal proxy; real liquid-response conclusions are based on RGB or liquid-level observations.

\Cref{tab:real_fixed_path_ablation} reports a representative fixed-path internal-ablation result.
The comparison includes a basic baseline, a smooth-only baseline, and a slosh-aware variant.
Each method is repeated three times.
All trials reach the goal, and the valid configurations use $v_{\mathrm{ref}}=0.20\,\mathrm{m/s}$, matched linear/angular velocity and acceleration limits, enabled fixed closed-loop delay compensation, disabled hard model constraint, and no tracking-protection stop.
Statistics are computed over the 10--90\% path-progress window.

\begin{table}[t]
  \centering
  \scriptsize
  \setlength{\tabcolsep}{3pt}
  \renewcommand{\arraystretch}{1.05}
  \caption{Representative fixed-path physical internal-ablation results over the 10--90\% path-progress window.}
  \label{tab:real_fixed_path_ablation}
  \label{tab:real_reduction_selected}
  \resizebox{\columnwidth}{!}{%
  \begin{tabular}{lccc}
    \toprule
    Metric & \Bzero{} & \Bsmooth{} & \Bslosh{} \\
    \midrule
    \multicolumn{4}{l}{\textit{A. Physical ablation statistics}}\\
    Model p95 [mm] & $1.100\pm0.034$ & $0.994\pm0.165$ & $0.632\pm0.101$ \\
    Model max [mm] & $2.102\pm0.222$ & $2.227\pm0.240$ & $1.255\pm0.039$ \\
    Model RMS [mm] & $0.547\pm0.016$ & $0.491\pm0.066$ & $0.307\pm0.035$ \\
    RGB p95 [mm] & $1.009\pm0.214$ & $1.272\pm0.283$ & $0.901\pm0.015$ \\
    RGB max [mm] & $2.398\pm0.482$ & $2.267\pm0.309$ & $2.162\pm0.405$ \\
    RGB RMS [mm] & $0.541\pm0.081$ & $0.639\pm0.060$ & $0.537\pm0.008$ \\
    Projection p95 [cm] & $2.947\pm0.087$ & $3.646\pm1.188$ & $3.442\pm0.128$ \\
    $|\omega|$ p95 [rad/s] & $0.234\pm0.015$ & $0.202\pm0.022$ & $0.242\pm0.033$ \\
    Arrival time [s] & $33.753\pm0.355$ & $35.169\pm0.537$ & $36.851\pm0.231$ \\
    \midrule
    \multicolumn{4}{l}{\textit{B. Relative reduction of \Bslosh{} with respect to \Bsmooth{}}}\\
    Model p95 & -- & -- & $36.4\%$ \\
    Model RMS & -- & -- & $37.4\%$ \\
    RGB p95 & -- & -- & $29.2\%$ \\
    RGB RMS & -- & -- & $15.9\%$ \\
    \bottomrule
  \end{tabular}}
\end{table}

The comparison between \Bsmooth{} and \Bzero{} shows that stronger smoothing can reduce $|\omega|$ p95 and some model-proxy statistics, but the RGB liquid envelope does not decrease accordingly.
\Bslosh{} does not have the lowest $|\omega|$ p95, yet both the model-proxy and RGB p95/RMS metrics are lower than those of \Bsmooth{}.
This indicates that the observed liquid-response improvement in these representative trials cannot be explained by smaller angular velocity alone.
Panel B of \cref{tab:real_fixed_path_ablation} reports the corresponding relative reductions.

Relative to the smooth-only baseline, \Bslosh{} reduces the model-proxy p95 and RMS by 36.4\% and 37.4\%, and reduces the RGB $H_{\mathrm{vis}}$ p95 and RMS by 29.2\% and 15.9\%.
Peak statistics are retained in \cref{tab:real_fixed_path_ablation}, but the isolated sensitivity of max values to visual envelope spikes makes p95/RMS the primary evidence in this comparison.
Relative to \Bzero{}, \Bslosh{} reduces the model proxy and RGB p95, while the RGB RMS is nearly unchanged; therefore \Bzero{} is used as a reference baseline rather than the main reduction denominator.
Overall, these representative fixed-path physical trials show consistent p95/RMS liquid-response improvement of the slosh-aware variant over the smooth-only baseline.
At the same time, \Bslosh{} has a longer arrival time, and its projection-error and angular-velocity statistics are not uniformly best.
The conclusion is therefore limited to liquid-response improvement rather than dominance across all task metrics.

The model-external comparison is used only as an interpretive diagnostic.
In this result, the internal model proxy $H_{\mathrm{model}}$ and the external RGB liquid envelope decrease consistently when comparing \Bslosh{} against \Bsmooth{}, supporting a trend-level explanation based on explicit slosh prediction.
The proxy still cannot replace external liquid observations, frame-level phase fidelity, ground-truth free-surface height, or hard safety monitoring.

\subsection{Online Computation and Execution Diagnostics}

Online solve performance is reported separately from liquid-response metrics to verify that \SPMPC{} remains a deployable receding-horizon local planner rather than an offline anti-slosh trajectory optimizer.
Relevant quantities include OCP solve-time distribution, control frequency, solver failures, timeouts, and extreme cases.
Median and high-percentile solve times are compared with the control period.

Execution diagnostics explain discrepancies in the physical system.
They include optimizer output, final base velocity command, command saturation, terminal deceleration, sensing delay, and base response.
These diagnostics support experimental interpretation but are not independent anti-slosh contributions.
When the execution chain substantially modifies the optimizer output, the corresponding physical conclusions are interpreted under that execution condition.
